\section*{Chapter \ref{ch:nw:the-north-american-map} --- The North American Map}

\begin{solution}{exo:nw:the-north-american-map:1}
$10^6/c_0 = \qty{3.336}{\milli\second}$ in vacuum and $1.4620$ times that, \qty{4.877}{\milli\second}, in fibre. Over the
\qty{1180.9}{\kilo\metre} between Aurora and Carteret the floors are \qty{3939.0}{\micro\second} and \qty{5758.8}{\micro\second}, a difference of
\qty{1.82}{\milli\second}: the prize that sent the fastest routes into the air.
\end{solution}

\begin{solution}{exo:nw:the-north-american-map:2}
Markham and Carteret are \qty{553.1}{\kilo\metre} apart: the one-way floor in vacuum is \qty{1.845}{\milli\second}, so a round trip cannot take
less than \qty{3.69}{\milli\second}. The \qty{2.05}{\milli\second} is therefore one-way, with a \omterm{def:nw:the-north-american-map:factor}{route factor} of 1.111 in air.
\end{solution}

\begin{solution}{exo:nw:the-north-american-map:3}
$2\,974.6 \times 1.0003 / c_0 = \qty{9.9}{\micro\second}$. The fastest published route between the two buildings is \qty{42}{\micro\second} above its
floor in air: a yardstick wrong by 10 of those 42 microseconds cannot rank routes that differ by a few.
\end{solution}

\begin{solution}{exo:nw:the-north-american-map:4}
The floor in air over \qty{1179.0}{\kilo\metre} is \qty{3933.8}{\micro\second}; $3\,986 / 3\,933.8 = 1.013$. The excess of
\qty{52.2}{\micro\second} is worth $52.2 \times 10^{-6} \times c_0 / 1.0003 = \qty{15.7}{\kilo\metre}$ of extra path in air (or the equivalent in
equipment).
\end{solution}

\begin{solution}{exo:nw:the-north-american-map:5}
Carteret--NY5 is \qty{25.96}{\kilo\metre}, \qty{126.6}{\micro\second} of straight fibre; Carteret--Mahwah \qty{55.6}{\kilo\metre},
\qty{271.0}{\micro\second}. The NY5 piece waits $271.0 - 126.6 = \qty{144.4}{\micro\second}$. The real figure is the difference between the
measured latencies of the firm's actual routes (longer than the geodesic, with switches, and different for each carrier), plus each venue's
own gateway time, so it is measured, not computed.
\end{solution}

\begin{solution}{exo:nw:the-north-american-map:6}
Markham: its building is sourced to TMX and ICE but its street address only to a directory. Three kilometres move its floors by up to
$3 \times 3.336 = \qty{10.0}{\micro\second}$ in vacuum and \qty{14.6}{\micro\second} in fibre, if the error lies along the route; an error across
it changes almost nothing.
\end{solution}

\begin{solution}{exo:nw:the-north-american-map:7}
\texttt{geodesic\_m(41.80, -88.24, 40.58, -74.25)} gives \qty{1180.13}{\kilo\metre}: \qty{1.13}{\kilo\metre} more than the paper's
\qty{1179}{\kilo\metre} and \qty{0.74}{\kilo\metre} less than the geodesic between the unrounded coordinates. Rounding to two decimals moves each
point by up to about \qty{0.6}{\kilo\metre}, which is \qty{2}{\micro\second} of light: fine for a paper, not for ranking routes.
\end{solution}

\begin{solution}{exo:nw:the-north-american-map:8}
The \omterm{def:nw:the-north-american-map:factor}{route factor} measures the fibre route against the floor \emph{in fibre}. Moving the same path into the air divides the glass part by
$1.4620/1.0003$, a saving of 31.6\% of the time spent in the medium, whatever the \omterm{def:nw:the-north-american-map:factor}{route factor}; the equipment (amplifiers against radios) and the
path (a cable's right of way against a tower line's) also change, so the new \omterm{def:nw:the-north-american-map:factor}{route factor} must be estimated separately. The statement confuses
the route's detour with the medium's speed.
\end{solution}

\begin{solution}{pb:nw:the-north-american-map:1}
\textbf{Part I.}
\begin{enumerate}
\item \qty{1180.87}{\kilo\metre} on WGS-84.
\item \qty{1177.90}{\kilo\metre}: \qty{2.97}{\kilo\metre} short, an error of $-0.25\%$.
\item \qty{1.68}{\kilo\metre}, \qty{5.6}{\micro\second} of light.
\item Every real route runs on or above the surface: cables in the ground, radios on towers. The chord is a physical bound nobody can use.
\end{enumerate}
\textbf{Part II.}
\begin{enumerate}[resume]
\item \qty{3939.0}{\micro\second}.
\item \qty{3940.1}{\micro\second}; between $n = 1.00027$ and $n = 1.0003$ the floor moves by \qty{0.12}{\micro\second}, nothing.
\item \qty{5758.8}{\micro\second}.
\item \qty{11.52}{\milli\second} in fibre, \qty{7.88}{\milli\second} in vacuum.
\end{enumerate}
\textbf{Part III.}
\begin{enumerate}[resume]
\item Cermak--Carteret is \qty{1129.2}{\kilo\metre}; the fibre floor is \qty{5506.9}{\micro\second}, and $6\,650/5\,506.9 = 1.208$.
\item \qty{1143}{\micro\second}, worth $1\,143 \times 10^{-6} \times c_0/1.4620 = \qty{234}{\kilo\metre}$ of glass: path, equipment and the ends.
\item The metro leg's floor is \qty{255.1}{\micro\second}; at the same \omterm{def:nw:the-north-american-map:factor}{route factor} it takes \qty{308.0}{\micro\second}, for
\qty{6.958}{\milli\second} in all, and $6\,958/5\,758.8 = 1.208$ against the Aurora--Carteret floor in fibre: going through Cermak, which
is almost on the way, costs little beyond the route's own factor.
\item $6\,550/5\,506.9 = 1.189$: \qty{100}{\micro\second} less, the equivalent of \qty{20.5}{\kilo\metre} of glass, from straighter segments or
faster equipment.
\end{enumerate}
\textbf{Part IV.}
\begin{enumerate}[resume]
\item \emph{Named result}: between Carteret and Aurora the one-way \omterm{def:nw:the-north-american-map:floor}{latency floor} is \qty{3.939}{\milli\second} in vacuum,
\qty{3.940}{\milli\second} in air and \qty{5.759}{\milli\second} in fibre along the geodesic; the 2010 fibre from Cermak ran at a \omterm{def:nw:the-north-american-map:factor}{route factor} of
1.21 over its own floor.
\item $5\,758.8 - 3\,982 = \qty{1777}{\micro\second}$ under the fibre floor: no fibre of standard glass can match it, however straight.
\item \qty{41.9}{\micro\second}, a \omterm{def:nw:the-north-american-map:factor}{route factor} of 1.011: about \qty{12.5}{\kilo\metre} of detour and equipment.
\item Only with a \omterm{def:nw:the-north-american-map:floor}{group index} below $3\,982/3\,939.0 = 1.011$ even along the geodesic: a fibre whose light travels almost as in air, the
hollow-core fibre of chapter~13, and a nearly straight path.
\item An error of \qty{0.3}{\kilo\metre} along the route is \qty{1}{\micro\second}; the geocodes place the buildings, not their \omterm{def:nw:colocation-products-and-how-they-are-sold:mmr}{meet-me rooms},
so the named result is good to about a microsecond.
\item It adds the Aurora--Cermak metro, at least \qty{255}{\micro\second} each way in fibre, to every use of the Cermak-based route for trading
Aurora's prices: a route built for one geography loses value when the venue moves.
\item Because firms trade from racks: the tails from the building's edge to the rack are part of the race, and only a rack-to-rack figure can
be compared across providers.
\item The geodesic gives the physical floor that no vendor can beat, so a vendor's latency can be read as a \omterm{def:nw:the-north-american-map:factor}{route factor} and compared with the
best that has been done.
\end{enumerate}
\end{solution}

\begin{solution}{iq:nw:the-north-american-map:1}
The NYSE group matches in Mahwah, Nasdaq in Carteret, Cboe's U.S. equities and options in Secaucus (NY5), all in New Jersey within about
\qty{56}{\kilo\metre} of each other; CME's futures in Aurora, Illinois, about \qty{1180}{\kilo\metre} away.
\iqlookfor{Named buildings, not just cities; the triangle's scale (tens of km, about 100--\qty{200}{\micro\second}) against Chicago's
(\qty{1200}{\kilo\metre}, about \qty{4}{\milli\second}); awareness that sites move and must be re-verified.}
\end{solution}

\begin{solution}{iq:nw:the-north-american-map:2}
About \qty{1180}{\kilo\metre} of geodesic: $1180/299\,792 = \qty{3.94}{\milli\second}$ in vacuum; times 1.462, \qty{5.76}{\milli\second} in fibre;
the 2010 fibre added about 20\% to its floor, the best published microwave route about 1\%.
\iqlookfor{The \qty{3.34}{\micro\second\per\kilo\metre} and \qty{4.88}{\micro\second\per\kilo\metre} rules; one-way against round trip; the
difference between a floor and a route.}
\end{solution}

\begin{solution}{iq:nw:the-north-american-map:3}
Light travels in air at nearly $c_0$ and in glass at $c_0/1.46$, so over the same path microwave is about 32\% faster. Fibre carries far more
bandwidth (the Toronto wireless service of this chapter sells 1 to 10 megabits a second; one fibre pair carries the 10- and 100-gigabit links of
chapter~1), and glass underground does not care about the weather, which a line of radios in the open air does (chapter~14).
\iqlookfor{\omterm{def:nw:the-north-american-map:floor}{Group index}; bandwidth and weather as the price of speed; firms using both, microwave for the few messages that race, fibre for
everything else (chapters~13 and~14).}
\end{solution}

\begin{solution}{iq:nw:the-north-american-map:4}
Take each building's address from the venue or the operator, geocode it with a cited geocoder, and compute the geodesic on WGS-84 with a
tested method (Vincenty's or Karney's); cross-check with an independent position. It must be accurate to a few hundred metres, a microsecond of
light, because routes compete on tens of microseconds and a spherical formula errs by 0.25\%.
\iqlookfor{Ellipsoid, not sphere; tests against published geodesics; sources for coordinates; accuracy stated in time.}
\end{solution}

\begin{solution}{iq:nw:the-north-american-map:5}
The direct feed is produced in the building where Nasdaq matches; the CTA plan's consolidated feed is produced in Mahwah, so a Nasdaq quote in
one of its stocks travels to Mahwah, is processed, and travels back: at least \qty{371}{\micro\second} of light in vacuum over the
\qty{55.6}{\kilo\metre}, more in fibre, plus the processing. The physics part cannot be engineered away; the processing part can.
\iqlookfor{Geography of the SIP against the direct feeds; round trip, not one way; the distinction between the floor and the processing time.}
\end{solution}

\begin{solution}{iq:nw:the-north-american-map:6}
Compute the floor between the exact buildings (\qty{3.94}{\milli\second} in air for Aurora--Carteret): \qty{3.99}{\milli\second} is a \omterm{def:nw:the-north-american-map:factor}{route
factor} of 1.013, plausible for a microwave route, impossible for fibre. Then ask whether the figure is one-way, radio to radio or rack to rack,
median or guaranteed, and how it is measured, and test it yourself with timestamped captures at both ends (chapters~4 and~5).
\iqlookfor{The floor as a sanity check; the definitions behind a published figure; measurement before purchase; the SLA figure against the
typical one (chapter~15).}
\end{solution}
